\begin{abstract}
Energy-matching diffusion samplers such as iDEM learn to sample from a Boltzmann density by regressing a score network onto a self-normalized importance sampling (SNIS) estimate of the noised score at noised replay-buffer samples. We show that the population optimum of these objectives does not depend on the buffer. Even a perfect sampler is therefore not a fixed point of training: the model is pulled towards a biased field, and the bias grows quickly with noise level and dimension. We observe that every training pair already holds an exact posterior sample, namely the clean buffer point that produced the noised input. \emph{Anchored Energy Matching} (AEM) adds this point as one extra particle in the SNIS estimate. Drawing on conditional SMC, we prove three properties. The anchored target is exactly unbiased for any number of proposals when the buffer is exact, so the true score becomes a fixed point. Its variance never exceeds that of single-sample score matching. When the buffer is wrong, its bias is a contracted version of the buffer's error, with an explicit constant. On a Gaussian mixture benchmark, anchoring reduces the score MSE by up to $475\times$ at high noise. In full sampler training it removes iDEM's systematic within-mode distortion: on the 32-dimensional many-well target, iDEM inflates the within-mode variance $4\times$, while AEM recovers it to within $5\%$ and lowers the energy Wasserstein distance by up to $29\times$, at a 1\% higher cost. The contraction constant also predicts the method's limitation, which we confirm. Where proposals cannot reach a mode, the anchored target follows the buffer, so mode discovery must come from an exploration phase.
\end{abstract}
